\documentclass[a4paper, serif, 10pt]{moderncv}

%% moderncv
% color and style
\moderncvcolor{black}
\moderncvstyle{banking}

%% page layout/margins
\usepackage[top=2.5cm, bottom=2.5cm, left=1.5cm, right=1.5cm]{geometry}

\nopagenumbers{}

%% Babel config for Dutch language
% ref:
% - https://latex3.github.io/babel/guides/locale-dutch.html
\usepackage[dutch]{babel}

%% fontspec
\usepackage{fontspec}

\IfFontExistsTF{Linux Libertine}{
  \setmainfont[Ligatures={TeX, Common}, Numbers=OldStyle]{Linux Libertine}
}{}
\IfFontExistsTF{Linux Libertine O}{
  \setmainfont[Ligatures={TeX, Common}, Numbers=OldStyle]{Linux Libertine O}
}{}

%% CTeX
% CJK support, used to show CJK characters in the resume
%
% - heading=false: disable ctex heading
% - punct=kaiming: use kaiming punctuations styles for CJK
% - scheme=plain: use plain scheme, do not override \normalsize font size
% - space=auto: space settings for CJK characters
%
% CJK fonts are resolved per locale language, e.g. Noto Serif/Sans CJK JP
% for Japanese resumes, with fallback to the next region if not installed
%
% ref:
% - http://ctan.mirrorcatalogs.com/language/chinese/ctex/ctex.pdf
\usepackage[UTF8, heading=false, punct=kaiming, scheme=plain, space=auto]{ctex}

\IfFontExistsTF{Noto Serif CJK SC}{
  \setCJKmainfont{Noto Serif CJK SC}
}{}
\IfFontExistsTF{Noto Sans CJK SC}{
  \setCJKsansfont{Noto Sans CJK SC}
}{}

%% line spacing
\usepackage{setspace}
\setstretch{1.125}

%% URL styling - use normal text instead of monospace
\urlstyle{same}

%% Patch \httplink and \httpslink to handle full URLs
% The original moderncv macros always prepend http:// or https:// to URLs.
% This causes issues when the URL already has a protocol (e.g., https://example.com)
% resulting in malformed URLs like https://https://example.com.
% This patch checks if the URL already contains "://" and uses it directly if so.
% Note: We use \str_set:Nx (x-expansion) to expand macros like \@homepage first.
\ExplSyntaxOn
\str_new:N \g_yamlresume_url_str

\RenewDocumentCommand{\httpslink}{O{}m}{
  \str_set:Nx \g_yamlresume_url_str {#2}
  \str_if_in:NnTF \g_yamlresume_url_str {://}
    { \url{#2} }
    { \url{https://#2} }
}

\RenewDocumentCommand{\httplink}{O{}m}{
  \str_set:Nx \g_yamlresume_url_str {#2}
  \str_if_in:NnTF \g_yamlresume_url_str {://}
    { \url{#2} }
    { \url{http://#2} }
}
\ExplSyntaxOff

\name{Sanne de Vries}{}
\title{QA Engineer met passie voor kwaliteit en testautomatisering}
\phone[mobile]{+31 6 12345678}
\email{sanne.devries@example.com}
\homepage{https://sannedevries.nl}

\address{Keizersgracht 123, Amsterdam, Noord-Holland, Nederland, 1015 CJ}{}{}

\extrainfo{{\small \faLinkedin}\ \href{https://www.linkedin.com/in/sannedevries}{@sannedevries} {} {} {} • {} {} {} 
{\small \faGithub}\ \href{https://github.com/sannedevries}{@sannedevries}}

\begin{document}

\maketitle

\section{Basis Informatie}

\cvline{}{\begin{itemize}
\item Ervaren QA Engineer met een sterke focus op testautomatisering en kwaliteitsborging binnen Agile teams
\item Expert in het opzetten van schaalbare testframeworks en het integreren van testen in CI/CD-pipelines
\item Analytisch en detailgericht, met ruime ervaring in API-testen, performance testing en testmanagement
\item Teamspeler die bruggen slaat tussen development, operations en productmanagement om kwaliteit continu te verbeteren
\end{itemize}}
\section{Opleidingen}

\cventry{sep 2014–jul 2018}
        {Bachelor, Informatica, Score: 7.8}
        {Technische Universiteit Delft}
        {\url{https://www.tudelft.nl}}
        {}
        {\begin{itemize}
\item Sterke basis gelegd in softwareontwikkeling, algoritmen en kwaliteitszorg
\item Praktijkervaring opgedaan met testmethodieken en geautomatiseerde tests
\item Samenwerkingsvaardigheden ontwikkeld via groepsprojecten en presentaties
\end{itemize}
\textbf{Cursussen}: Datastructuren en Algoritmen, Software Engineering, Databases, Netwerken, Besturingssystemen, Testen van Software, Projectmanagement, Webtechnologie}
\section{Werkervaring}

\cventry{mrt 2022–Heden}
        {Senior QA Engineer}
        {ING}
        {\url{https://www.ing.com}}
        {}
        {\begin{itemize}
\item Verantwoordelijk voor de kwaliteitsstrategie van de mobiele bankieren-app en online bankieren
\item Opgezet en onderhouden van een schaalbaar testautomatiseringsframework met Cypress en TypeScript
\item Geïntegreerd van geautomatiseerde tests in CI/CD-pipelines, wat leidde tot 40\% snellere releases
\item Begeleiden van junior QA-engineers en het verzorgen van kennissessies over testautomatisering
\item Samenwerken met productowners en ontwikkelaars om acceptatiecriteria en testscenario's te definiëren
\end{itemize}
\textbf{Trefwoorden}: Testautomatisering, Selenium, Cypress, API-testen, CI/CD, Agile, Scrum, Kwaliteitszorg}

\cventry{jun 2019–feb 2022}
        {QA Engineer}
        {Adyen}
        {\url{https://www.adyen.com}}
        {}
        {\begin{itemize}
\item Ontworpen en uitgevoerd van teststrategieën voor de checkout-omgeving en betaaloplossingen
\item Geautomatiseerd van regressietests met Postman, Java en Python, resulterend in 60\% minder handmatig testwerk
\item Uitgevoerd van performance tests met JMeter en Gatling om schaalbaarheid en betrouwbaarheid te waarborgen
\item Actief bijgedragen aan Agile ceremonies en continue verbetering van het testproces
\item Gerapporteerd van kwaliteitsrisico's en testresultaten aan stakeholders
\end{itemize}
\textbf{Trefwoorden}: Teststrategie, Postman, Java, Python, Prestatietesten, Regressietesten}

\cventry{sep 2018–mei 2019}
        {Junior QA Engineer}
        {Coolblue}
        {\url{https://www.coolblue.nl}}
        {}
        {\begin{itemize}
\item Uitgevoerd van handmatige en geautomatiseerde testcases voor de webshop en mobiele app
\item Vastleggen en opvolgen van bugs in JIRA, in nauwe samenwerking met ontwikkelaars
\item Deelgenomen aan Scrum-sprints en bijgedragen aan het verbeteren van de Definition of Done
\item Opgesteld van testdocumentatie en testrapportages
\end{itemize}
\textbf{Trefwoorden}: Handmatig testen, Testcases, Bugtracking, JIRA, Scrum}
\section{Talen}

\cvline{Nederlands}{Moedertaal of tweetalige vaardigheid}
\cvline{Engels}{Volledige professionele vaardigheid \hfill \textbf{Trefwoorden}: Cambridge C1, IELTS 7.5}
\cvline{Duits}{Beperkte werkvaardigheid \hfill \textbf{Trefwoorden}: Goethe B1}
\section{Vaardigheden}

\cvline{Testautomatisering}{Expert \hfill \textbf{Trefwoorden}: Selenium, Cypress, Playwright, Cucumber, Robot Framework}
\cvline{API-testen}{Geavanceerd \hfill \textbf{Trefwoorden}: Postman, REST Assured, Swagger, GraphQL, SoapUI}
\cvline{Prestatietesten}{Intermediair \hfill \textbf{Trefwoorden}: JMeter, Gatling, LoadRunner, K6}
\cvline{Testmanagement}{Geavanceerd \hfill \textbf{Trefwoorden}: Teststrategie, Testplanning, Risicobeheer, ISTQB, TMMi}
\cvline{Programmeertalen}{Geavanceerd \hfill \textbf{Trefwoorden}: Java, Python, JavaScript, SQL, Bash}
\cvline{DevOps}{Intermediair \hfill \textbf{Trefwoorden}: Jenkins, GitLab CI, Docker, Kubernetes, Azure DevOps}
\section{Onderscheidingen}

\cventry{nov 2023}
        {ING}
        {QA Innovatie Award}
        {}
        {}
        {Uitgereikt voor het ontwikkelen van een innovatief testautomatiseringsframework dat de release-snelheid en kwaliteit van de mobiele bankieren-app aanzienlijk verbeterde.}
\section{Certificaten}

\cventry{mrt 2021}
        {ISTQB}
        {ISTQB Certified Tester Advanced Level}
        {\url{https://www.istqb.org}}
        {}
        {}

\cventry{sep 2020}
        {Scrum Alliance}
        {Certified ScrumMaster}
        {\url{https://www.scrumalliance.org}}
        {}
        {}
\section{Publicaties}

\cventry{jan 2023}
        {Testautomatisering in Agile teams: van strategie naar uitvoering}
        {ING Tech Blog}
        {\url{https://www.ing.com/techblog/testautomatisering}}
        {}
        {\begin{itemize}
\item Beschrijft een praktische aanpak voor het opzetten van testautomatisering binnen Agile teams
\item Deelt lessen uit de praktijk over tooling, samenwerking en continue verbetering
\item Geeft concrete handvatten voor het integreren van tests in CI/CD-pipelines
\end{itemize}}
\section{Projecten}

\cventry{jan 2022–dec 2022}
        {Een schaalbaar testautomatiseringsframework voor de mobiele bankieren-app.}
        {Testautomatiseringsframework voor ING Bankieren}
        {\url{https://github.com/ing/test-framework}}
        {}
        {\begin{itemize}
\item Ontworpen van een modulair framework met Cypress en TypeScript voor end-to-end tests
\item Geïntegreerd van tests in de CI/CD-pipeline met automatische rapportage
\item Gerealiseerd van een onderhoudbare testsuite die de regressietijd met 50\% verkortte
\end{itemize}
\textbf{Trefwoorden}: Cypress, TypeScript, CI/CD, Appium, API-testen}

\cventry{mrt 2020–nov 2021}
        {Een geautomatiseerde performance teststraat voor de checkout-omgeving.}
        {Performance teststraat voor Adyen Checkout}
        {\url{https://github.com/adyen/performance-suite}}
        {}
        {\begin{itemize}
\item Opgezet van gedistribueerde load tests met JMeter en Gatling
\item Geautomatiseerd van testuitvoering en rapportage in Kubernetes
\item Gevisualiseerd van resultaten in Grafana voor snelle analyse en bijsturing
\end{itemize}
\textbf{Trefwoorden}: JMeter, Gatling, Kubernetes, Grafana}
\section{Interesses}

\cvline{Hardlopen}{Marathon, Trailrunning, Intervaltraining}
\cvline{Fotografie}{Portretfotografie, Landschapsfotografie, Lightroom}
\cvline{Reizen}{Azië, Europa, Wandelen}
\section{Vrijwilligerswerk}

\cventry{sep 2019–jun 2023}
        {Vrijwilliger Programmeerles}
        {CodeUur}
        {\url{https://www.codeuur.nl}}
        {}
        {\begin{itemize}
\item Verzorgd van programmeerlessen aan basisschoolkinderen om hen kennis te laten maken met technologie
\item Begeleid van leerlingen bij het bouwen van eenvoudige games en websites
\item Bijgedragen aan het enthousiasmeren van kinderen voor STEM-onderwijs
\end{itemize}}
    
\end{document}